% code_and_data.tex -- Appendix A
\label{sec:code_and_data}

This appendix describes the companion repository
\nolinkurl{github.com/JackDMenendez/dcl-paper-02-sm-derivation}
(v1.0 tag), with the aim that any reader can audit or reproduce
any claim in the paper from the repository contents.  The
audit-table convention of Section~\ref{subsec:audit_convention}
applies: every claim has a row in Table~\ref{tab:audit} with a
status tag and a pointer to its supporting artifact (a script, a
derivation, or a measured value cited from the literature).

\subsection{Repository layout}
\label{subsec:repo_layout}

The repository is organised as follows.

\begin{itemize}
\item \nolinkurl{paper/} --- LaTeX source for this paper.
  \begin{itemize}
  \item \nolinkurl{paper/main.tex} --- the document root.
  \item \nolinkurl{paper/sections/} --- one file per section,
    including \nolinkurl{audit_table.tex} (the canonical
    record of every PASS / PART / STUB / FAIL claim status).
  \item \nolinkurl{paper/macros/} --- shared LaTeX packages and
    commands.
  \item \nolinkurl{paper/paper-bib/references.bib} --- the
    bibliography source.
  \end{itemize}

\item \nolinkurl{src/utilities/} --- sympy verification scripts.
  Thirteen scripts in total: six inherited from Paper~I (the
  \nolinkurl{automorphism_*.py} and \nolinkurl{tick_rule_*.py}
  files) and seven new to Paper~II that close Phases~1--4 of
  \nolinkurl{notes/work_plan.md}.  Each script is self-contained
  and runs from a Python virtual environment with \texttt{sympy}
  (see Appendix~\ref{sec:reproducibility} for setup).  The output
  of each run is the operational PASS/PART/FAIL for the
  corresponding audit-table row; running every script is the
  end-to-end verification of Paper~II's claims.

\item \nolinkurl{notes/} --- working theoretical notes that the
  paper either cites (via the audit table's evidence column) or
  builds on.  At v1.0 the directory contains:
  \begin{itemize}
  \item \nolinkurl{work_plan.md} --- the phased plan for Phases~0
    through 5, with resolution stanzas for each closed phase.
  \item Per-phase findings and structural interpretations
    (one note per closed Phase or sub-phase):
    \begin{itemize}
    \item \nolinkurl{su3_branch_consistency.md}
    \item \nolinkurl{chirality_alignment.md}
    \item \nolinkurl{cp_modification_obstruction.md}
    \item \nolinkurl{su3_generation_from_colour_memory.md}
    \item \nolinkurl{aut_centralizer_enumeration.md}
    \item \nolinkurl{aut_centralizer_extras_commutators.md}
    \item \nolinkurl{induced_gauge_action_nonabelian.md}
    \end{itemize}
  \item Structural-synthesis and framing notes:
    \begin{itemize}
    \item \nolinkurl{mass_chirality_coupling.md}
    \item \nolinkurl{debt_to_measurement.md}
    \item \nolinkurl{no_spacetime_torsion.md}
    \end{itemize}
  \item \nolinkurl{extras_as_a1_accounting.md} (\texttt{DRAFT})
    --- the hypothesis that the 59 extras of the Phase~3
    centralizer are the algebraic bookkeeping channels for
    $\mathcal{A}=1$ during what an SM observer identifies as
    decoherence.  Sketches the follow-on script
    \nolinkurl{decoherence_as_projection.py}, deferred from v1.0.
  \item \nolinkurl{lie_algebra_proof_sketch.md} --- a pointer to
    Paper~I's working proof sketch
    (\nolinkurl{external/dcl/notes/lie_algebra_automorphism_proof_sketch.md},
    accessible via the Windows directory junction described in
    \nolinkurl{CLAUDE.md}).
  \end{itemize}

\item \nolinkurl{release_notes/} --- per-version change logs and
  GitHub Release bodies (see also
  Appendix~\ref{subsec:repro_release_flow}).

\item \nolinkurl{CLAUDE.md} --- working brief: project status,
  what NOT to change, cross-references to Paper~I and the parallel
  formalisation effort in physics-research, notes index.  This
  file is the entry point for anyone (human or AI) returning to
  the project.

\item \nolinkurl{CITATION.cff} --- machine-readable citation
  metadata.

\item \nolinkurl{external/} --- Windows directory junctions to
  Paper~I (\nolinkurl{external/dcl}) and the physics-research
  formalisation effort (\nolinkurl{external/research}), gitignored
  per \nolinkurl{.gitignore}.  On a fresh clone these need to be
  recreated locally to follow cross-references to those
  repositories; the seeded prose explicitly cites Paper~I's
  Zenodo DOI~\cite{menendez2026geometry} so that a reviewer
  without those local checkouts can still navigate the published
  artifact.
\end{itemize}

\subsection{The audit table as canonical record}
\label{subsec:audit_canonical}

The convention introduced in Section~\ref{subsec:audit_convention}
is the load-bearing one for code/data audit: prose claims defer
to Table~\ref{tab:audit}, every row has a status tag, every
status tag has a pointer to a verifiable artifact, and the table
itself is part of the released artifact (frozen at v1.0).  A
reader who doubts any specific claim in the paper can navigate to
the relevant audit-table row, follow the evidence-column pointer
to the supporting script, and re-run the script for independent
verification.  Negative results (\texttt{FAIL} rows --- Phase~1.5
in particular) are part of the contribution and are recorded with
the same convention.

The current v1.0 audit-table state is:

\begin{itemize}
\item Established factors (\texttt{PASS}, $8$ rows): discrete
  spatial $\mathrm{Aut}(\Gamma, \mathbf{V}) = O_h$ of order 48;
  RGB sublattice contributes only $\mathbb{Z}_3 \subset SU(3)$;
  $SO(3,1) \times U(1)$ acts on the existing $\mathbb{C}^2$
  (dim 7); direct-product structure on extended $\mathbb{C}^{12}$
  (dim 18); tick-rule consistency on $\mathbb{C}^{12}$; bipartite
  Wilson plaquette gauge invariance; SU(3) representation branch
  consistency; non-abelian $SU(3)$ from $\mathbb{C}^3$ memory.
\item Partial / characterisation (\texttt{PART}, $3$ rows):
  SM-chirality coupling alignment, exact equality vs containment
  in $\mathrm{Aut}_\text{ext}$, explicit $1/g^2$ prefactor for
  the non-abelian Wilson actions.
\item Disconfirmed (\texttt{FAIL}, $1$ row): modified tick rule
  for Branch~B SU(3) (no antilinear modification admits SM-style
  $CP$).
\item Central claim (\texttt{PART}): Eq.~(137) of Paper~I --
  containment $\supseteq$ holds; equality $=$ does not; the
  characterisation outcome is the framework's terminal
  prediction.
\end{itemize}

\subsection{Provenance and license}
\label{subsec:provenance}

Paper text and figures are released under the Creative Commons
Attribution 4.0 (CC BY 4.0) licence.  Source code in the
\nolinkurl{src/utilities/} directory is released under the MIT
licence.  The six inherited scripts (\nolinkurl{automorphism_*.py}
and \nolinkurl{tick_rule_*.py}) are verbatim copies of the v1.0
versions in Paper~I~\cite{menendez2026geometry}; edits to those
scripts in Paper~II are additive (extensions of the verification,
not refactorings of the existing calculation) so that Paper~I's
printed output remains the canonical reference for the original
six rows.

The complete commit history is available at the repository's
\nolinkurl{main} branch; the v1.0 release tag pins the exact
commit corresponding to this paper's published artifact.  Issues
and pull requests are accepted at the GitHub repository.
